\chapter{Diseases of the Spleen}
\label{ch:spleen}
The spleen is the largest lymphoid organ in the body, located in the left upper quadrant of the abdomen. It functions as a blood filter, removing old and damaged red blood cells (culling), and plays important roles in immune function, including IgM production, complement activation, and opsonization. The red pulp is responsible for filtering blood, while the white pulp contains lymphoid tissue involved in immune responses. Splenomegaly (enlarged spleen) and hypersplenism are common manifestations of various underlying conditions.

%-------------------------------------------------------------
\section{Asplenia and Functional Asplenia}
%-------------------------------------------------------------

%-------------------------------------------------------------

%-------------------------------------------------------------

%-------------------------------------------------------------

%-------------------------------------------------------------

%-------------------------------------------------------------

\label{sec:Asplenia and Functional Asplenia}
%-------------------------------------------------------------%-------------------------------------------------------------

\subsection{Asplenia}
\label{sec:Asplenia}
Asplenia is the absence of splenic function, which may be congenital (asplenia syndrome, often associated with complex congenital heart disease and situs abnormalities) or acquired (splenectomy, sickle cell disease, celiac disease, inflammatory bowel disease). The condition results from anatomical absence or functional impairment of the spleen, which normally clears encapsulated bacteria through opsonin-dependent phagocytosis. Clinical features include increased risk for overwhelming post-splenectomy infection (OPSI), a fulminant sepsis with 50--70\% mortality. Diagnosis is based on history of splenectomy or conditions causing functional asplenia, confirmed by the presence of Howell-Jolly bodies on blood smear or absence of splenic uptake on nuclear medicine scan. Management involves pneumococcal, meningococcal, and Hib vaccination, prophylactic antibiotics (penicillin V or amoxicillin), and patient education about infection risks. Prognosis is good with preventive measures; OPSI carries high mortality. Functional asplenia (impaired splenic function without anatomical absence) occurs in sickle cell disease (autosplenectomy from repeated splenic infarction), celiac disease, and inflammatory bowel disease.

%-------------------------------------------------------------
\section{Hypersplenism}
%-------------------------------------------------------------

%-------------------------------------------------------------

%-------------------------------------------------------------

%-------------------------------------------------------------

%-------------------------------------------------------------

%-------------------------------------------------------------

\label{sec:section:Hypersplenism}
%-------------------------------------------------------------%-------------------------------------------------------------

\subsection{Hypersplenism}
\label{sec:Hypersplenism}
Hypersplenism is a syndrome of splenic overactivity resulting in cytopenias (one or more of anemia, leukopenia, thrombocytopenia) due to excessive sequestration and destruction of blood cells. It is always secondary to an underlying condition causing splenomegaly. Clinical features include cytopenias with corresponding symptoms (anemia, infection risk, bleeding tendency). Diagnosis is suggested by cytopenias, splenomegaly, and bone marrow showing adequate or hypercellular hematopoiesis. Management is directed at the underlying cause; splenectomy may be considered for severe, symptomatic cytopenias refractory to medical management, though it increases the risk of overwhelming post-splenectomy infection (OPSI). Prognosis depends on the underlying condition and severity of cytopenias.

%-------------------------------------------------------------
\section{Splenic Enlargement (Splenomegaly)}
%-------------------------------------------------------------

%-------------------------------------------------------------

%-------------------------------------------------------------

%-------------------------------------------------------------

%-------------------------------------------------------------

%-------------------------------------------------------------

\label{sec:Splenic Enlargement (Splenomegaly)}
%-------------------------------------------------------------%-------------------------------------------------------------

\subsection{Massive Splenomegaly}
\label{sec:Massive Splenomegaly}
Massive splenomegaly is spleen enlargement extending below the umbilicus or crossing the midline. Common causes in developing countries include malaria, visceral leishmaniasis (kala-azar), schistosomiasis, and myelofibrosis; in Western countries, myeloproliferative neoplasms (chronic myeloid leukemia, polycythemia vera, primary myelofibrosis), lymphomas (hairy cell leukemia, splenic marginal zone lymphoma), and storage diseases (Gaucher disease) are more common. The condition results from the underlying infiltrative, congestive, or neoplastic process. Clinical features include significant symptoms and complications including splenic infarction, portal hypertension, and severe cytopenias from hypersplenism. Diagnosis is based on imaging and identification of the underlying cause. Management is directed at the underlying cause; splenectomy may be required for symptomatic relief or life-threatening complications. Prognosis depends on the underlying etiology.

\subsection{Splenomegaly}
\label{sec:Splenomegaly}
Splenomegaly is enlargement of the spleen beyond its normal size (normally approximately 12 cm in length). It has numerous causes classified as congestive (portal hypertension, heart failure, hepatic vein thrombosis), infectious (malaria, infectious mononucleosis, infective endocarditis, visceral leishmaniasis, tuberculosis), hematologic (hemolytic anemias, myeloproliferative disorders, lymphomas, leukemias, myelofibrosis), inflammatory/autoimmune (sarcoidosis, Felty syndrome, SLE, rheumatoid arthritis), and infiltrative (Gaucher disease, Niemann-Pick disease, amyloidosis). Clinical features depend on the degree of enlargement; mild splenomegaly may be asymptomatic, while massive splenomegaly causes early satiety, left upper quadrant pain, and discomfort, with a firm, smooth spleen on palpation. Diagnosis involves identifying the underlying cause through blood tests, imaging (ultrasound, CT), and bone marrow biopsy when indicated. Management is directed at the underlying condition. Prognosis depends on the underlying etiology.

%-------------------------------------------------------------
\section{Splenic Infarction}
%-------------------------------------------------------------

%-------------------------------------------------------------

%-------------------------------------------------------------

%-------------------------------------------------------------

%-------------------------------------------------------------

%-------------------------------------------------------------

\label{sec:section:Splenic Infarction}
%-------------------------------------------------------------%-------------------------------------------------------------

\subsection{Splenic Infarction}
\label{sec:Splenic Infarction}
Splenic infarction is ischemic necrosis of splenic parenchyma resulting from occlusion of the splenic artery or its branches. The most common cause is embolization from the heart (atrial fibrillation, endocarditis, mural thrombus), followed by hematologic disorders (myeloproliferative disorders, sickle cell disease, thalassemia) and vasculitis. Clinical features include acute left upper quadrant pain, sometimes radiating to the left shoulder (Kehr's sign). Diagnosis is based on CT scan showing wedge-shaped, hypodense, non-enhancing lesions. Management is conservative (analgesics, observation) with treatment of the underlying cause. Prognosis is generally favorable; rarely, complicated infarction (abscess, rupture) may require splenectomy.

%-------------------------------------------------------------
\section{Splenic Neoplasms}
%-------------------------------------------------------------

%-------------------------------------------------------------

Splenic neoplasms comprise primary benign lesions, primary hematolymphoid malignancies (such as splenic marginal zone lymphoma), and secondary metastatic tumors. Because of the spleen's rich vascular structure and phagocytic role, neovascular or lymphoproliferative growth alters splenic filtering function and produces splenomegaly. Evaluation includes abdominal cross-sectional imaging, peripheral blood and bone marrow analysis, and therapeutic or diagnostic splenectomy when indicated.

%-------------------------------------------------------------

%-------------------------------------------------------------

%-------------------------------------------------------------

%-------------------------------------------------------------

\label{sec:Splenic Neoplasms}
%-------------------------------------------------------------%-------------------------------------------------------------

\subsection{Splenic Cysts}
\label{sec:Splenic Cysts}
Splenic cysts are fluid-filled lesions in the spleen, classified as true cysts (with epithelial lining, including congenital, epidermoid, and dermoid cysts) or pseudocysts (without epithelial lining, usually resulting from prior trauma, infarction, or infection). Most are asymptomatic and discovered incidentally on imaging. Clinical features of large cysts include abdominal pain, early satiety, or a palpable mass. Diagnosis is based on imaging. Management is observation for small, asymptomatic cysts; surgical excision (cystectomy or partial splenectomy) or splenectomy for large or symptomatic cysts; percutaneous aspiration and sclerotherapy is an alternative for selected cases. Prognosis is excellent.

\subsection{Splenic Lymphoma}
\label{sec:Splenic Lymphoma}
Splenic lymphoma is involvement of the spleen by systemic lymphoma or, less commonly, primary splenic origin. Splenic marginal zone lymphoma (SMZL) is the most common primary splenic lymphoma, a low-grade B-cell lymphoma. The condition results from lymphoproliferation; hepatitis C virus infection is associated with SMZL in some cases. Hairy cell leukemia is another important splenic lymphoproliferative disorder characterized by the BRAF V600E mutation. Clinical features include massive splenomegaly, lymphocytosis (villous lymphocytes on blood smear in SMZL), and pancytopenia (in hairy cell leukemia). Diagnosis involves bone marrow biopsy, flow cytometry, and CT imaging. Management varies by histology: splenectomy and/or rituximab for SMZL; cladribine or moxetumomab for hairy cell leukemia. Prognosis depends on histologic subtype and disease stage.

%-------------------------------------------------------------
\section{Splenic Rupture and Trauma}
%-------------------------------------------------------------

%-------------------------------------------------------------

%-------------------------------------------------------------

%-------------------------------------------------------------

%-------------------------------------------------------------

%-------------------------------------------------------------

\label{sec:Splenic Rupture and Trauma}
%-------------------------------------------------------------%-------------------------------------------------------------

\subsection{Splenic Rupture}
\label{sec:Splenic Rupture}
Splenic rupture is a life-threatening emergency, most commonly caused by blunt abdominal trauma (motor vehicle accidents, falls, sports injuries). The spleen is the most commonly injured abdominal organ in trauma. The condition may also occur spontaneously in the setting of massive splenomegaly (infectious mononucleosis, hematologic malignancies, myelofibrosis). Clinical features include left upper quadrant pain, peritoneal signs, hemodynamic instability (tachycardia, hypotension), and Kehr's sign (left shoulder pain). Diagnosis is based on focused assessment with sonography for trauma (FAST) and CT scan, which also grades the injury. Management follows Advanced Trauma Life Support (ATLS) principles: hemodynamic resuscitation, splenic grading, and decision between non-operative management (for hemodynamically stable patients with lower-grade injuries) and emergent splenectomy (for hemodynamic instability or high-grade injuries). Prognosis depends on injury severity and timely intervention; post-splenectomy vaccination (pneumococcal, meningococcal, \textit{Haemophilus influenzae} type b) and prophylactic antibiotics are essential to prevent OPSI.

